%% Interstitial: Paper 1 -> Paper 2
%% Philosophical and conceptual bridge: prediction to dynamics.

Chapter~\ref{chap:paper-one} isolated the predictive content of a single
query at a single time.  The minimal predictive query $Q_*$ is the
sufficient statistic for prediction; the predictive kernel encodes the
conditional structure of $P$ relative to $Q$.  The observer's present
measurement has a definite predictive face --- one that is forced, not
chosen.

\medskip

\paragraph{The residual freedom.}
What Paper~1 does not determine is how these predictive structures relate
\emph{across time}.  The observer makes measurements at time $0$ and
time $t$; each has a minimal predictive query and a predictive kernel.
But the relationship between them is not, after Paper~1, settled.  The
system evolves; the observer's predictive capacity evolves with it.
Nothing in the construction of $Q_*$ at time $0$ constrains the
predictive structure at time $t > 0$.

The freedom is this: compatible observable laws force $P$, and $P$
forces $Q_*$ at each time, but the temporal relationship between
$Q_*|_{t=0}$ and $Q_*|_{t=s}$ appears unconstrained.  Two observers
with the same $P$ but different ideas about how predictive content
propagates through time are not ruled out by anything established
so far.

\medskip

\paragraph{The elimination.}
Chapter~\ref{chap:paper-two} shows this temporal freedom is not free.
The \emph{predictive operator family} $\{K_t\}_{t \geq 0}$ organises
predictive content across time --- $K_t g$ is the predictive operator
applied to the function $g$ evaluated $t$ time units ahead --- and
the central result is the \emph{semigroup property}:
\[
  K_{t+s} = K_t \circ K_s.
\]
This is not an assumption about the dynamics.  It is a theorem about
the structure of prediction under compatible observable laws: temporal
coherence of prediction is forced by the same observable compatibility
that forced $P$.  The observer's present and future are not two
independent predictions stitched together by choice; they are aspects
of a single coherent structure.

Chapter~\ref{chap:paper-two} also establishes \emph{Koopman-Perron
duality}: the predictive operator $K_t$ and the Koopman operator are
dual under the $L^2(P)$ inner product.  The observable-layer dynamics
and the latent state-space dynamics are not competitors; they are dual
aspects of the same structure, neither primary.  The program does not
eliminate the state space; it shows that neither layer is the ground.

\medskip

\paragraph{The structure this introduces.}
The semigroup $\{K_t\}$ is abstract: it holds for any observable
dynamics satisfying the program's hypotheses, without specifying the
form of the queries or the geometric structure of the outcome spaces.
The observer knows that their predictive capacity is temporally
coherent --- but in what form?  For a real observer embedded in time,
with bounded and local discriminative capacity, the semigroup must be
instantiated concretely.  That is the question Chapter~\ref{chap:paper-three}
takes up.
